Actualmente, el control de drones se realiza principalmente mediante controles remotos tradicionales, los cuales pueden representar un desafío para usuarios con poca experiencia. El manejo de estos dispositivos requiere un proceso de aprendizaje que implica coordinación entre la percepción visual del entorno y la ejecución precisa de comandos de control, lo que demanda un alto nivel de atención y habilidad por parte del operador. En muchos casos, esto hace necesario recibir capacitación previa o incluso contar con certificaciones específicas para operar drones de forma segura.

Como resultado, estas barreras técnicas pueden limitar la accesibilidad de esta tecnología para nuevos usuarios o para aplicaciones donde se busca una interacción más intuitiva y natural con el sistema. Esto reduce significativamente el acceso del público general al uso de drones, principalmente debido a la pronunciada curva de aprendizaje que enfrentan los nuevos operadores. Esta situación genera un impacto adicional, ya que la escasez de pilotos capacitados puede limitar el uso de drones en diversas áreas de aplicación, como la industria, operaciones de bomberos, monitoreo ambiental, reconocimiento de zonas y otras tareas donde estas plataformas podrían aportar beneficios significativos.

En este contexto, el presente trabajo propone el desarrollo de un sistema de control de drones que permita reducir las barreras de acceso asociadas a los métodos tradicionales de pilotaje. Para ello, se plantea la implementación de una interfaz de control basada en movimientos corporales capturados mediante un sistema de seguimiento óptico, con el objetivo de ofrecer una forma de interacción más natural e intuitiva para el usuario. Mediante este enfoque se busca disminuir la curva de aprendizaje requerida para operar drones, facilitando su uso por personas con poca o ninguna experiencia previa. De esta manera, el sistema propuesto pretende contribuir al desarrollo de nuevas formas de interacción humano-drone que amplíen el acceso a esta tecnología y potencien su aplicación en entornos educativos, industriales y de investigación. Tomando esto en cuenta el presente trabajo puede impulsar el desarrollo de una alternativa de controlar drones más intuitiva para el ser humano y sin la necesidad de tener que tener experiencia.
